\section{The Frame of the Life: Chronology and the Gallio Anchor}

\subsection{What the letters and Acts yield by themselves}

Read alone, the sources yield only relative time. The letters give three
numbers---``after three years,'' ``after fourteen years'' (Gal 1:18; 2:1),
and ``above fourteen years ago'' for the ecstasy of 2~Corinthians 12:2---
plus one datable name, the Nabataean king Aretas (2~Cor 11:32). Acts gives
sequence and seasons, two-year blocks at Caesarea and Rome (24:27; 28:30),
a stay of ``a year and six months'' at Corinth (18:11), and a handful of
datable officials and events: the edict by which ``Claudius had commanded
all Jews to depart from Rome'' (18:2; Suetonius confirms the expulsion of
Jews making disturbances ``at the instigation of Chrestus,'' \work{Claudius}
25, but gives no year---the conventional 49 rests on the fifth-century
Orosius, not verified here), the
procuratorial succession ``Felix had for successor Portius Festus''
(24:27, datable only within a debated range, c.~55--60), and one proconsul
of Achaia: ``But when Gallio was proconsul of Achaia, the Jews with one
accord rose up against Paul, and brought him to the judgment seat'' (Acts
18:12). Without external control, none of this attaches to a calendar.

\subsection{The Delphi inscription}

The control exists for Gallio alone. Lucius Junius Gallio---brother of the
philosopher Seneca---is named in a fragmentary Greek inscription found at
Delphi and published early in the twentieth century: a letter of the
emperor Claudius to the city, in which Claudius bears the title
``acclaimed imperator the 26th time.'' The classic epigraphic study is
Adolf Deissmann's appendix ``The Proconsulate of L.~Junius Gallio''
(\work{St.~Paul: A Study in Social and Religious History}, English
translation 1912, pp.~235--260, with a facsimile of the stone), which this
study read in full in the public scan identified in the References.
Deissmann's argument, resting on Hermann Dessau's collation of Claudius's
titulature, runs: the 26th imperatorial acclamation belongs between
(at earliest, the end of 51) 25 January 52 and 1 August 52, when the 27th
is already attested; the letter names Gallio as proconsul of Achaia in
office; proconsuls served one year, entering office in early summer;
therefore Gallio's proconsulship ran from the summer of 51 to the summer
of 52 (Deissmann: he ``entered on his proconsulship in the summer
(nominally 1 July) of 51 A.D.''), with 52--53 the only alternative the
evidence tolerates. Since Acts places the hearing before Gallio at the end
of Paul's eighteen-month Corinthian ministry, Paul's arrival in Corinth
falls c.~49--50 and his departure c.~51--52. This is the one absolute
date in Paul's life, and the whole modern chronology hangs from it.
The Church's own teaching office has used the same anchor: ``It is the
most reliable date in the whole of his biography\ldots\ It says that
Gallio was Proconsul in Corinth between the years 51 and 53. Thus we have
one absolutely certain date'' (Benedict~XVI, General Audience, 27 August
2008).

\witnesslimit{Ceilings, stated plainly: this study verified Deissmann's
1912 edition and argument at first hand in a public-domain scan, and the
anchor as reported by a modern official witness; it did not collate the
stone, later-joined fragments, or the current critical edition of the
inscription. The inscription dates Gallio's tenure, not Paul's hearing
within it; ``summer 51--summer 52'' for the proconsulship and ``50--52''
for Paul's Corinth are the responsibly stated ranges, and the year-wide
uncertainty propagates through every derived date below.}

\subsection{From the anchor to a life}

Counting backward from Corinth c.~50: the Macedonian mission (Philippi,
Thessalonica, Beroea, Athens) fills the immediately preceding year or
two; the Jerusalem council, which the sequence of both Acts and Galatians
puts before the Aegean mission, lands c.~48--49; the ``fourteen years''
of Galatians 2:1 then reach back from it to a revelation c.~33--36, and
the ``three years'' of the Arabia--Damascus period nest inside that span.
Counting forward: Corinth to c.~52; Ephesus ``until Pentecost'' seasons
and a roughly three-year residence (Acts 19:8--10; 20:31) to c.~52--55/57;
the last journey to Jerusalem with the collection c.~57 (Pentecost, Acts
20:16); two years' custody at Caesarea under Felix to c.~57--59; the
voyage, wintering at Malta, and arrival in Rome c.~60; ``two whole years''
of open custody to c.~62 (Acts 28:30); death under Nero between 64 and
68. Every one of these figures except the anchor is a reconstruction
(class K/S), and honest rivals exist at nearly every joint: the
crucifixion in 30 or 33; the council as early as 47 or as late as 51;
Festus's accession anywhere from 55 to 60, dragging the Caesarea--Rome
block with it; and a minority of modern reconstructions redistributes the
Acts journeys around the letters' own hints far more radically. The
detailed table in the terminal appendix states each event with its basis
and its status; the narrative that follows uses the conventional frame
just derived, always as a frame and never as data.

\begin{fourstable}{Joint}{Conventional}{Alternatives}{What moves it}
\properrow{Crucifixion of Jesus}{30 or 33}{both defended}{astronomy of
Nisan 14; every early Pauline date shifts with it}
\properrow{Revelation near Damascus}{c.~33--36}{as early as 31}{whether
Gal 1--2's 3 $+$ 14 years are consecutive or concurrent}
\properrow{Jerusalem council}{c.~48--49}{c.~47 to 51}{identification of
Gal 2 with Acts 15 or Acts 11; length of the first journey}
\properrow{Corinth (first stay)}{c.~50--52}{49--51 or 51--53}{Gallio's
year (51--52 vs.\ 52--53); where the hearing fell within it}
\properrow{Festus succeeds Felix}{c.~59}{55--60}{thin Roman evidence;
drags Caesarea, the voyage, and the Roman custody with it}
\properrow{Death}{67--68 (chronicle tradition)}{64 (persecution's first
wave)}{substance of the release/second-arrest tradition}
\end{fourstable}
